\chapter{Embedding the Engine in Java}
\label{ch:embedding}

The engine is a Java library. A program creates a \fn{Rete}, declares its
classes, loads rule files, asserts objects and fires; the shell and the
GUI are thin layers over exactly this API. This chapter shows the calls
that matter, then the optional modules.

\section{Class path}

An embedding needs \file{morendo-core-\version.jar} and the two Log4j
jars, \file{log4j-api} and \file{log4j-core}, from \file{libs/} of the
distribution; \file{jspecify} is optional (it carries the nullness
annotations). Add \file{morendo-examples} for the sample beans,
\file{morendo-gui} if rules should be able to call \fn{(view)}, and
\file{morendo-messaging} or \file{morendo-service} for the modules
described at the end of the chapter. Any \fn{FunctionGroup} listed in a
jar's \file{META-INF/services/org.morendo.rete.FunctionGroup} is
registered automatically when a \fn{Rete} is created.

\section{A complete example}

The following program was compiled and run against the \fn{examples}
module; it loads the rule file from Section~\ref{ch:java-objects} and
asserts two \fn{Account} beans.

\begin{lstlisting}[style=java]
import org.morendo.rete.Rete;
import woolfel.examples.model.Account;
import woolfel.examples.model.AccountHobby;
import woolfel.examples.model.Hobby;

import java.io.PrintWriter;
import java.util.List;

public class Example {
    public static void main(String[] args) throws Exception {
        Rete engine = new Rete();
        try {
            engine.addPrintWriter("stdout", new PrintWriter(System.out, true));
            engine.declareObject(Account.class);
            engine.declareObject(AccountHobby.class);
            engine.declareObject(Hobby.class);
            engine.loadRuleset("samples/ruleset/sample1.clp");

            Account ann = new Account();
            ann.setFirst("ann");
            ann.setLast("lee");
            ann.setAge(34);
            Account bob = new Account();
            bob.setFirst("bob");
            bob.setLast("ray");
            bob.setAge(12);
            engine.assertObjects(List.of(ann, bob));

            int fired = engine.fire();
            System.out.println(fired + " rule(s) fired");

            bob.setAge(31);   // the bean notifies the engine, which re-matches
            System.out.println(engine.fire() + " rule(s) fired after the change");
        } finally {
            engine.close();
        }
    }
}
\end{lstlisting}

\begin{lstlisting}[style=shell]
woolfel.examples.model.Account
woolfel.examples.model.AccountHobby
woolfel.examples.model.Hobby
ann lee is over 30
1 rule(s) fired
bob ray is over 30
1 rule(s) fired after the change
\end{lstlisting}

The three class names are printed by \fn{declareObject}. The second
\fn{fire} shows the property-change path: \fn{setAge} on the
\fn{Account} bean fires a \fn{PropertyChangeEvent}, the engine (a
registered listener) modifies the shadow fact, and the \fn{Over\_30} rule
matches Bob.

\section{The Rete API}

\subsection{Life cycle and threading}

\begin{lstlisting}[style=java]
Rete engine = new Rete();
...
engine.addCloseHook(() -> pool.release(engine));
engine.close();
boolean done = engine.isClosed();
\end{lstlisting}

Creating an engine starts a daemon thread for its message router;
\fn{close()} clears everything, stops that thread and runs the close
hooks. An engine is \emph{not thread-safe}: drive it from one thread, or
post commands through the message router (below). The \fn{(exit)}
function is \fn{close()}; nothing in the engine calls \fn{System.exit}.

\subsection{Serving many threads}
\label{sec:threads}

Nothing inside an engine is synchronized: its working memory, agenda,
scopes and profiling counters belong to the thread that uses it, and
two threads asserting into one engine corrupt it. A program that serves
requests from many threads therefore gives each thread its own engine,
or keeps a pool and hands an engine to one thread at a time. What is
shared between engines in a JVM is small: the strategy registry of
\fn{defstrategy}, and in the messaging module the content-handler and
agent registries; templates, rules, facts, functions and the network
are per engine. The cost of an engine is its compiled network and one
daemon thread for its message router, plus whatever facts it holds, so
a pool of a dozen engines is unremarkable.

The pattern is the one the service module implements: load the rules
once per engine, and for each request check an engine out, assert the
request's facts, fire, read the results and check it back in after
retracting what the request added, so the next request starts from the
same state. \fn{Rete.clearFacts()} and \fn{clearObjects()} empty an
engine completely when a request should leave nothing behind, rules
and templates included in neither. The message router is the other way
in: several threads may post commands to one engine through their own
channels, and the router's command thread runs them one at a time
(Section~\ref{sec:router}).

\subsection{Loading rules}

\begin{lstlisting}[style=java]
engine.loadRuleset("rules/orders.clp");            // file path, URL or classpath:...
engine.loadRuleset(inputStream);
engine.build("(defrule r (order (total ?t)) => (printout t ?t crlf))");
\end{lstlisting}

\fn{loadRuleset} is the \fn{batch} function: the location may be a file
path relative to the working directory, an \fn{http:}, \fn{file:} or
\fn{jar:} URL, or \fn{classpath:dir/file}. Parse errors are written to
the engine output rather than thrown, so register a print writer before
loading if you want to see them. \fn{build} executes CLIPS text held in
a string.

\subsection{Templates and objects}

\begin{lstlisting}[style=java]
engine.declareObject(Order.class);                       // template "com.acme.Order"
engine.declareObject(Order.class, "Order");              // template "Order"
engine.declareObject(RushOrder.class, "RushOrder", "Order");   // with a parent template

engine.assertObject(order, null, false, true);    // dynamic shadow fact, template from the class
engine.assertObjects(List.of(a, b, c));            // the same for each element
engine.modifyObject(order);                        // after changing a bean that fires no events
engine.retractObject(order);
Fact shadow = engine.getShadowFact(order);         // null if not asserted
\end{lstlisting}

\fn{assertObject} takes the object, an optional template name (null
means the class's own template, and the fact is also propagated to
parent templates), a \fn{static} flag and a \fn{shadow} flag. A static
object is reference data: it gets no change listener,
\fn{modifyObject} ignores it, and \fn{resetObjects} re-asserts it. With
\fn{shadow} true the property values are copied into the fact; with it
false the fact reads the live object on every access, with no copy and
no listener, which is cheaper for large objects that change rarely.
Asserting an instance that is already asserted does nothing.

Plain facts are built from a template:

\begin{lstlisting}[style=java]
Deftemplate order = (Deftemplate) engine.findTemplate("order");
Fact f = order.createFact(new Object[] {
        new Slot("id", 42), new Slot("customer", "acme")}, -1);
engine.assertFact(f);
engine.retractFact(f);             // or engine.retractById(f.getFactId())
\end{lstlisting}

\subsection{Running and inspecting}

\begin{lstlisting}[style=java]
int fired = engine.fire();            // drain the agenda of the module in focus
engine.fire(1);                       // one activation
List<Rule> rules = engine.getRulesFired();   // the distinct rules fired so far
List<Fact> facts = engine.getAllFacts();
List<Object> objects = engine.getObjects();
engine.setWatch(Rete.Watch.RULES);    // FACTS, ACTIVATIONS, ALL
engine.setProfile(Rete.Profile.ALL);
engine.setFocus("MAIN");
Object limit = engine.getDefglobalValue("limit");
engine.declareDefglobal("limit", 18);
\end{lstlisting}

\subsection{Output}

\begin{lstlisting}[style=java]
StringWriter out = new StringWriter();
engine.addPrintWriter("capture", out);
...
PrintWriter w = engine.removePrintWriter("capture");   // the caller closes it
\end{lstlisting}

Every writer registered with \fn{addPrintWriter} is a console: it
receives everything the engine prints to \fn{t}, which is \fn{printout
t} output, watch traces and error messages. This is how the golden tests
capture a run and how a server can attach a log. Output to another router
name goes only to the router opened under that name, by \fn{(open)} in
CLIPS or from Java:

\begin{lstlisting}[style=java]
engine.openRouter("audit", new FileWriter("audit.log", true)); // or (file, mode)
engine.writeMessage("checked\n", "audit");   // what (printout audit ...) does
engine.closeRouter("audit");                  // closeRouter(null) closes all
engine.setInputSupplier(() -> nextLine());    // where (readline t) reads from
\end{lstlisting}

Without an input supplier, \fn{readline t} reads standard input.
\fn{Rete.close()} closes every open router.

\subsection{Listeners}

\begin{lstlisting}[style=java]
engine.getRuleCompiler().addListener(new CompilerListener() {
    public void ruleAdded(CompileEvent e)    { log.info(e.getMessage()); }
    public void ruleRemoved(CompileEvent e)  { }
    public void compileError(CompileEvent e) { log.warn(e.getMessage()); }
});
\end{lstlisting}

\fn{CompilerListener} is called when rules are added, removed, or fail
to compile (a missing template, an unknown function, a parse error);
\fn{CompileEvent.getRule()} and \fn{getMessage()} say which and why.
\fn{addEngineEventListener} registers an \fn{EngineEventListener} that
receives an \fn{EngineEvent} of kind \fn{ASSERT} or \fn{RETRACT}, with
the fact, every time working memory changes; listeners run on the thread
that asserted, so keep them short.

\subsection{Temporal facts}

\begin{lstlisting}[style=java]
engine.assertTemporalObject(event, null, Instant.now(),
        Instant.now().plus(Duration.ofMinutes(5)), false);
engine.assertFact(temporalFact, effective, expiration);   // expiration may be null
\end{lstlisting}

The times are stored on the shadow fact as epoch milliseconds
(Chapter~\ref{ch:temporal}).

\subsection{Functions written in Java}

A function implements the \fn{Function} interface of
\fn{org.morendo.rete}: \fn{getName()}, \fn{getReturnType()},
\fn{executeFunction(Rete, Parameter[])} returning a \fn{ReturnVector},
\fn{getParameter()} and \fn{toPPString()} for the usage text. Read arguments with \fn{params[i].getValue(engine,
ValueType.OBJECT)}, which resolves bound variables and nested calls.

\begin{lstlisting}[style=java]
public final class Shout implements Function {
    public String getName() { return "shout"; }
    public ValueType getReturnType() { return ValueType.STRING; }
    public Class<?>[] getParameter() { return new Class<?>[] {ValueParam.class}; }
    public String toPPString(Parameter[] params, int indents) { return "(shout <string>)"; }
    public ReturnVector executeFunction(Rete engine, Parameter[] params) {
        String s = String.valueOf(params[0].getValue(engine, ValueType.STRING));
        DefaultReturnVector rv = new DefaultReturnVector();
        rv.addReturnValue(new DefaultReturnValue(ValueType.STRING, s.toUpperCase()));
        return rv;
    }
}
\end{lstlisting}

Register it from Java or from CLIPS:

\begin{lstlisting}[style=java]
engine.declareFunction(new Shout());
\end{lstlisting}
\begin{lstlisting}
(load-function com.acme.Shout)
\end{lstlisting}

Several functions are bundled in a \fn{FunctionGroup}: \fn{getName()}
names the group, \fn{loadFunctions(Rete)} declares each function and
records it, and \fn{listFunctions()} returns that list, which is what
\fn{(functions AcmeFunctions)} prints.

\begin{lstlisting}[style=java]
public final class AcmeFunctions implements FunctionGroup {
    private final List<Function> functions = new ArrayList<>();

    public String getName() { return "AcmeFunctions"; }

    public void loadFunctions(Rete engine) {
        for (Function f : List.of(new Shout())) {
            engine.declareFunction(f);
            functions.add(f);
        }
    }

    public List<Function> listFunctions() { return functions; }
}
\end{lstlisting}

\fn{(load-function-group com.acme.AcmeFunctions)} loads the group from
CLIPS. To have every engine pick it up from the class path instead, list
the class in the jar's service file (\file{META-INF/services/}, file
\file{org.morendo.rete.FunctionGroup}). The example beans' module has
another small function, \fn{HelloFunction}.

A few habits of the built-in functions are worth copying. Name a function
with lower-case words separated by hyphens, \fn{str-trim} rather than
\fn{strTrim}. Check the number and kind of the arguments before using
them, and throw an \fn{IllegalArgumentException} with a clear message
when they are wrong; the shell prints it, and a rule that passes the wrong
thing is then easy to find. Keep each function to one task and keep no
state in it: a value that must survive between calls belongs in a
\fn{defglobal}, which each engine keeps separately, and state shared by
several engines needs a thread-safe class of its own
(Section~\ref{sec:threads}). A common way to work is to prototype a
function as a \fn{deffunction} at the shell, and rewrite it in Java once
it is settled and known to matter for speed.

\subsection{Driving the engine through the message router}
\label{sec:router}

To accept CLIPS text from another thread, or to get the engine's output
and results as events, use a channel on the message router, as the shell
does:

\begin{lstlisting}[style=java]
MessageRouter router = engine.getMessageRouter();
StringChannel channel = router.openChannel("MyApp");    // InterestType.MINE
router.setCurrentChannelId(channel.getChannelId());     // (printout t ...) goes here

channel.executeCommand("(batch samples/only/only_1.clp)", true);  // true: wait
List<MessageEvent> events = new ArrayList<>();
channel.fillEventList(events);
for (MessageEvent e : events) {
    switch (e.getType()) {
        case ENGINE -> System.out.print(e.getMessage());      // printed output
        case RESULT -> System.out.println(e.getMessage());    // the expression's value
        case ERROR  -> ((Exception) e.getMessage()).printStackTrace();
        default -> { }
    }
}
engine.close();
\end{lstlisting}

The router's command thread parses and executes each expression, so
several channels may post commands concurrently and the engine still
runs them one at a time. \fn{InterestType.ALL} makes a channel receive
the output of every channel, which is what the GUI's log tab uses.

\subsection{Logging}

Log4j~2 is the only logging dependency. The default configuration in
the core jar writes WARN and above to standard error; set
\fn{-Dmorendo.log.level=DEBUG} for detail, or supply your own
configuration with \fn{-Dlog4j2.configurationFile}. The alternative
\fn{log4j2-file.xml} in the same jar also writes
\file{logs/morendo.log}, with \fn{-Dmorendo.log.file} to move it.

\section{The service module}

\fn{morendo-service} wraps engines for server use: a
\fn{RuleService} holds one or more \emph{rule applications}, each with a
pool of engines initialized from a JSON configuration, and hands out an
\fn{EngineContext} per request:

\begin{lstlisting}[style=java]
EngineContext ctx = service.getEngine("test", null);   // waits for a free engine
try {
    ctx.assertObject(order, false, true);
    ctx.executeRules();
} finally {
    ctx.close();      // retracts the request's facts, returns the engine
}
\end{lstlisting}

The point of the pool is that creating a \fn{Rete} is quick, but
declaring the model classes, compiling the rules and asserting the
initial data can take far longer, and none of it should be repeated for
each request. The pool does that work once per engine; the request code
only asserts its data, fires and reads the results.

Each application has an \fn{EnginePool}, safe to use from any number of
threads. \fn{getEngine} hands out an idle engine, creates one while the
pool is below \fn{maxPool}, and otherwise waits up to the application's
\fn{checkoutTimeout} (milliseconds, 5000 by default) for one to come
back, returning null when none does. \fn{close()} retracts the facts and
objects the request added, so the initial data loaded from the
configuration stays and the next request starts from it; call
\fn{keepFacts()} on the context first to leave them in place.
\fn{service.getEnginePool(name, version)} exposes the pool, and the
administration object's \fn{reload} operations act on the engines that
are idle at the time.

The configuration (\file{samples/configuration/sample\_config.json})
names, per application, the model classes to declare, the function
groups to load, the rule files, and the initial data as fact files, JSON
or serialized objects; the pool sizes \fn{initialPool} and
\fn{maxPool} and the \fn{checkoutTimeout}; and a service name. In a servlet
container, register \fn{org.morendo.service.servlet.RuleStartupService}
as a listener in \file{web.xml}, as \file{samples/configuration/web.xml}
does:

\begin{lstlisting}[style=shell]
<listener>
  <listener-class>org.morendo.service.servlet.RuleStartupService</listener-class>
</listener>
\end{lstlisting}

Put the configuration in \file{/WEB-INF/Ruleconfig.json}, with a capital
R, and the \Morendo\ core and service jars, Log4j and Jackson in
\file{/WEB-INF/lib}. When the web application starts, the listener reads
the configuration, builds the pools and stores itself, a
\fn{RuleService}, as a servlet-context attribute named after the
configuration's \fn{serviceName}. A servlet looks it up by that name:

\begin{lstlisting}[style=java]
RuleService service = (RuleService) getServletContext().getAttribute("sample");
EngineContext ctx = service.getEngine("test", null);
try {
    ...                     // assert the request's data
    ctx.executeRules();
} finally {
    ctx.close();
}
\end{lstlisting}

The service keeps simple statistics over all its engines, updated each
time a context is closed: \fn{getRequests()},
\fn{getAverageResponseTime()}, \fn{getAverageRulesFired()} and
\fn{getTotalRulesFired()}.

Outside a container, \fn{RuleServiceImpl.createInstance(file)} reads
the same configuration, builds the applications and returns the service;
call \fn{initialize()} to create the engine pools. An application may
carry a \fn{version} in the configuration, and \fn{getEngine} takes the
name and version (null when none was given). Several versions of one
application can therefore run side by side in one service. That matters
where old rules stay in force for old business: an insurer rating a
back-dated change to a policy must use the rules in effect on that date,
and may want to compare the result with the current version's, without a
second server. Initial data of the
\fn{jsonData} kind is deserialized into the class named by its
\fn{name} field; \fn{clipsData} is loaded from its URL with
\fn{load-facts}.

\section{The messaging module}

\fn{morendo-messaging} connects an engine to a Jakarta Messaging topic.
The JMS provider's client jars and a running broker are the user's; the
module carries only the API. From CLIPS:

\begin{lstlisting}
(create-message-client "org.apache.activemq.jndi.ActiveMQInitialContextFactory"
    "tcp://localhost:61616" "ConnectionFactory" "dynamicTopics/rules"
    "user" "password" "bus")
(register-content-handler org.morendo.messaging.TextHandler)
(send-msg "bus" "(assert (order (id 7)))")
(close-message-client "bus")
\end{lstlisting}

The client is stored as a global under the name given last. Incoming
messages are dispatched by their JMS type to registered
\fn{ContentHandler}s; the built-in \fn{TextHandler} treats a text
message as CLIPS and executes it in the engine, so two engines on the
same topic can send each other facts and commands. Subscriptions are
non-durable and messages published by the client itself are not
delivered back to it. The \emph{agent} functions keep a registry of
remote engines (\fn{register-agent}, \fn{ping-agent},
\fn{pprint-agents}) built on the same client.
